%! Sinusoidal Function Context and Data Modeling — Unit 3, Lesson 3.7 — Group Activity
\documentclass[10pt]{article}
\usepackage{precalculus-article}
\usepackage{precalculus-boxes}
\newcommand{\TallMath}[1]{$\displaystyle #1\rule[-1.4em]{0pt}{3.2em}$}

\begin{document}

\pageheader{Unit 3, Lesson 3.7}{Group Activity}
\namepartnerperiod

\vspace{0.10in}

%----------------------------------------------------------------------
% Context
%----------------------------------------------------------------------
\begin{scenariobox}[Context: Modeling Periodic Real-World Data]{sky}
Sinusoidal functions appear in temperature records, tidal data, daylight hours,
and many other natural phenomena. In this activity your group will build and
evaluate a sinusoidal model from real-world data. Choose \textbf{one tier}
based on your readiness. Everyone should be ready to explain the reasoning
behind each parameter your group chooses.
\end{scenariobox}

\vspace{0.08in}

%----------------------------------------------------------------------
% Tier R
%----------------------------------------------------------------------
\begin{tcolorbox}[colback=white, colframe=black!40,
    title=\textbf{Tier R --- Reinforce},
    fonttitle=\bfseries, arc=2mm, left=3mm, right=3mm, top=2mm, bottom=2mm]

\textbf{Context:} The table below gives the average monthly temperature (${}^\circ$F)
in a small coastal town. The maximum and minimum are already labeled for you.

\vspace{4pt}
\begin{center}
\renewcommand{\arraystretch}{1.35}
\small
\begin{tabular}{|c|c|c|c|c|c|c|c|c|c|c|c|c|}
\hline
Month $x$ & 1 & 2 & 3 & 4 & 5 & 6 & 7 & 8 & 9 & 10 & 11 & 12 \\
\hline
Temp $T$ (${}^\circ$F) & 48 & 51 & 58 & 65 & 72 & 78 & \textbf{80} & 79 & 72 & 63 & 54 & \textbf{46} \\
\hline
\end{tabular}
\end{center}

\vspace{6pt}
\textbf{Step 1 --- Amplitude and Midline.}
The labeled max is \textbf{80} (month 7) and min is \textbf{46} (month 12 / month 1 area).

$A = \dfrac{80 - 46}{2} = $ \underline{\hspace{2cm}} \qquad
$D = \dfrac{80 + 46}{2} = $ \underline{\hspace{2cm}}

\vspace{6pt}
\textbf{Step 2 --- Period and $B$.}
The maximum occurs at month 7. The next maximum would occur at month
\underline{\hspace{1.2cm}} (one year later), so the period $k = $ \underline{\hspace{1.2cm}}.

$B = \dfrac{2\pi}{k} = $ \underline{\hspace{3cm}}

\vspace{6pt}
\textbf{Step 3 --- Phase Shift.}
Using a cosine model, $C = x_{\max} = $ \underline{\hspace{1.5cm}}.

\vspace{6pt}
\textbf{Step 4 --- Write the Model.}
$T(x) = A\cos(B(x - C)) + D = $ \underline{\hspace{7cm}}

\vspace{6pt}
\textbf{Step 5 --- Validate.}
Predict $T(1)$ using your model: $T(1) \approx$ \underline{\hspace{2.5cm}}.
Actual value from the table: 48. How close is your prediction? \writeline

\vspace{6pt}
\textbf{Step 6 --- Predict.}
Use your model to estimate the average temperature in month 4.5 (mid-May).
$T(4.5) \approx $ \underline{\hspace{2.5cm}} ${}^\circ$F.

\end{tcolorbox}

\vspace{0.08in}

%----------------------------------------------------------------------
% Tier A
%----------------------------------------------------------------------
\begin{tcolorbox}[colback=white, colframe=black!40,
    title=\textbf{Tier A --- Apply},
    fonttitle=\bfseries, arc=2mm, left=3mm, right=3mm, top=2mm, bottom=2mm]

\textbf{Context:} The table below shows the height of the tide (in feet above
mean sea level) recorded every hour at a harbor over one tidal cycle.

\vspace{4pt}
\begin{center}
\renewcommand{\arraystretch}{1.35}
\small
\begin{tabular}{|c|c|c|c|c|c|c|c|c|c|c|c|c|c|}
\hline
$t$ (hr) & 0 & 1 & 2 & 3 & 4 & 5 & 6 & 7 & 8 & 9 & 10 & 11 & 12 \\
\hline
$h(t)$ (ft) & 1.5 & 3.8 & 5.5 & 6.2 & 5.5 & 3.8 & 1.5 & $-0.8$ & $-2.5$ & $-3.2$ & $-2.5$ & $-0.8$ & 1.5 \\
\hline
\end{tabular}
\end{center}

\vspace{6pt}
\textbf{1.}\quad Identify the maximum and minimum values and the times at which they occur.

Max $= $ \underline{\hspace{1.5cm}} at $t = $ \underline{\hspace{1.2cm}} hr.
\quad Min $= $ \underline{\hspace{1.5cm}} at $t = $ \underline{\hspace{1.2cm}} hr.

\vspace{6pt}
\textbf{2.}\quad Compute $A$ and $D$.

$A = \dfrac{\max - \min}{2} = $ \underline{\hspace{3cm}} \qquad
$D = \dfrac{\max + \min}{2} = $ \underline{\hspace{3cm}}

What does the midline represent in the context of tidal data? \writeline

\vspace{6pt}
\textbf{3.}\quad Estimate the period $k$ from the data. Explain how you found it.
\writelines{2}

$k = $ \underline{\hspace{1.5cm}} hr \quad $\Rightarrow$ \quad $B = \dfrac{2\pi}{k} = $ \underline{\hspace{3cm}}

\vspace{6pt}
\textbf{4.}\quad Using a cosine model with $C = t_{\max}$, write the full sinusoidal model.

$h(t) = $ \underline{\hspace{8cm}}

\vspace{6pt}
\textbf{5.}\quad Validate: compute $h(0)$ from your model and compare to the actual value.

$h(0) = $ \underline{\hspace{3cm}} \quad Actual: 1.5 ft \quad Residual: \underline{\hspace{2cm}}

\vspace{6pt}
\textbf{6.}\quad A ship requires a minimum tide height of 2 ft to enter the harbor safely.
Using your model, at approximately what times during the 12-hour period is entry safe?
\writelines{3}

\end{tcolorbox}

\vspace{0.08in}

%----------------------------------------------------------------------
% Tier E
%----------------------------------------------------------------------
\begin{tcolorbox}[colback=white, colframe=black!40,
    title=\textbf{Tier E --- Extend},
    fonttitle=\bfseries, arc=2mm, left=3mm, right=3mm, top=2mm, bottom=2mm]

\textbf{Context:} The tables below give average monthly high temperatures for two
cities, City A (coastal) and City B (inland continental).

\vspace{4pt}
\textbf{City A:}
\begin{center}
\renewcommand{\arraystretch}{1.3}
\small
\begin{tabular}{|c|c|c|c|c|c|c|c|c|c|c|c|c|}
\hline
Month & 1 & 2 & 3 & 4 & 5 & 6 & 7 & 8 & 9 & 10 & 11 & 12 \\
\hline
Temp (${}^\circ$F) & 58 & 60 & 63 & 66 & 69 & 72 & 74 & 75 & 73 & 68 & 63 & 59 \\
\hline
\end{tabular}
\end{center}

\vspace{4pt}
\textbf{City B:}
\begin{center}
\renewcommand{\arraystretch}{1.3}
\small
\begin{tabular}{|c|c|c|c|c|c|c|c|c|c|c|c|c|}
\hline
Month & 1 & 2 & 3 & 4 & 5 & 6 & 7 & 8 & 9 & 10 & 11 & 12 \\
\hline
Temp (${}^\circ$F) & 28 & 34 & 50 & 66 & 78 & 88 & 94 & 90 & 76 & 60 & 42 & 31 \\
\hline
\end{tabular}
\end{center}

\vspace{6pt}
\textbf{1.}\quad For each city, compute $A$, $D$, $k$, $B$, and $C$ (using cosine), then
write the sinusoidal model.
\writelines{5}

\vspace{4pt}
\textbf{2.}\quad Compare your models. Which city has the larger amplitude? What does this
mean physically?
\writelines{2}

\vspace{4pt}
\textbf{3.}\quad For each city, compute the residual at $x = 1$ (January). Which model fits
January better?
\writelines{3}

\vspace{4pt}
\textbf{4.}\quad Compute residuals at $x = 4, 7, 10$ for City B. Based on the residuals,
do you think a sinusoidal model is a good fit for City B? Explain.
\writelines{3}

\vspace{4pt}
\textbf{5.}\quad State the contextual domain for each model. Using City B's model, predict
the average high temperature in month 14. Is this prediction trustworthy? Explain.
\writelines{3}

\end{tcolorbox}

\end{document}
